% !PRINTED 607
\weberpage{607}

\begin{center}
\textbf{\S\,164.\qquad \gesperrt{The Class Field.}}
\end{center}

To the number group $A$ in the quadratic field $\Omega$ there shall exist an algebraic field over $\Omega$, to be denoted by $\mathfrak{K}$ or $\mathfrak{K}(A)$, of which we shall assume only:

7. \gesperrt{Definition of the class field.}

\gesperrt{The prime ideals $\mathfrak{p}_1$ of the first degree of the principal class $\overline{A}_1$, and only these, shall split again in the field $\mathfrak{K}(A)$ into prime ideals of the first degree.}

Thus it is only the prime ideals in $\Omega$, which we designated earlier by $\mathfrak{p}_1$, that split in $\mathfrak{K}$ into prime ideals of the first degree.

Any finite number of primes may thereby be excepted. They are then included in the excludent.

If $n$ is the relative degree of the field $\mathfrak{K}$ over $\Omega$, and $p$ a natural prime which the ideals $\mathfrak{p}_1$ and $\mathfrak{P}$ in $\Omega$ and $\mathfrak{K}$ divide, and if one denotes the absolute norms in the fields $\Omega$ and $\mathfrak{K}$ by $N_\Omega, N_\mathfrak{K}$, then
\[
N_\mathfrak{K}(\mathfrak{P}) = N_\Omega(\mathfrak{p}_1) = p, \quad N_\mathfrak{K}(\mathfrak{p}_1) = p^n,
\]
and consequently $\mathfrak{p}_1$ must split in the field $\mathfrak{K}$ into $n$ prime ideals $\mathfrak{P}$:
\weq{1}{\mathfrak{p}_1 = \mathfrak{P}_1 \mathfrak{P}_2 \ldots \mathfrak{P}_n.}
The prime ideals which divide the fundamental numbers of the fields $\Omega$ or $\mathfrak{K}$ shall always be excluded. Then the $\mathfrak{P}_1, \mathfrak{P}_2, \ldots, \mathfrak{P}_n$ are distinct from one another.

Whether such a field $\mathfrak{K}$ exists remains for the present undecided. We will investigate what can be concluded from the definition.

Since $n$ ideals $\mathfrak{P}$ divide every ideal $\mathfrak{p}_1$, we have
\weq{2}{\prod\nolimits \left[1 - N_\mathfrak{K}(\mathfrak{p}_1)^{-s}\right]^n = \prod\nolimits \left[1 - N_\mathfrak{K}(\mathfrak{P})^{-s}\right],}
in which the first product extends over all ideals $\mathfrak{p}_1$ of the field $\Omega$, the second over all ideals $\mathfrak{P}$ of the first degree of the field $\mathfrak{K}$; and by Theorem I, § 197 of Band II, which holds for any field whatever,
\weq{3}{(s-1)\prod\nolimits \frac{1}{\left[1-N_\mathfrak{K}(\mathfrak{P})^{-s}\right]}}
has for $s=1$ a finite limiting value different from zero, which is there reduced to the class number in the field $\mathfrak{K}$. If we put, for brevity,
%CONT>>
